\documentstyle[11pt]{article}
\setlength{\textwidth}{170mm}
\setlength{\textheight}{230mm}
\setlength{\oddsidemargin}{0mm}
\setlength{\evensidemargin}{10mm}
\setlength{\topmargin}{-10mm}
\def\RR{\hbox{\raise.4ex\hbox{$\scriptstyle{|}$}\hskip-1.85pt R}}
\pagestyle{empty}
\begin{document}

\noindent
{\bf University of Illinois} \hfill \large{\bf Econ 476} \hfill {\bf Department of Economics} \\
{\bf Spring 2001} \hfill \large{{\bf Final Exam Review}}\hfill {\bf Roger Koenker}

\vspace{10mm}
%The final exam for 476 will be {\bf May 13 at 1:30 in Room 126 Commerce West.} 
%It will be closed note, closed book.

\vspace{5mm}
This handout includes a sampling of recent questions from 476 final exams.  
It should provide a reasonable sample to gauge the nature of the questions and 
to review the material covered in the course.  It could be expected that two exam questions will come from these review questions.


\begin{center}
*****
\end{center}

\vspace{10mm}
\noindent
{\bf University of Illinois} \hfill \large{\bf Econ 476} \hfill {\bf Department
of Economics} \\
{\bf Spring 1998} \hfill {\bf FINAL EXAMINATION}\hfill {\bf Roger Koenker}

\vspace{5mm}

\begin{enumerate}
%Q1
\item Let $Z$ be a $U[0, 1]$ random variable and define (hint:  draw pictures!)
\begin{quote}
$X = Z - \frac{1}{2}$\\
$Y=|2Z-1|-\frac{1}{2}$\\
$U= \mbox{ sgn } (X)$\\
$V= \mbox{ sgn } (Y)$
\end{quote}
\begin{enumerate}
\item Find the covariance matrix of the random vector $(X,Y, U,V).$
\item Which pairs of these variables, if any, are stochastically independent?
\end{enumerate}

%Q2
\item An important recent paper on semiparametric estimation of panel data
models begins with the observations that if two random variables $X$ and $Y$
are iid with common density $f(\cdot)$, then $Z=Z-Y$ has a symmetric
distribution centered at zero.  Explain.
\end{enumerate}



\vspace{10mm}
Spring 1997 \hfill {\bf Econ 476 Final Exam} \hfill Roger Koenker
\begin{enumerate}
%Q1.
\item Consider the problem of estimating the location parameter, $\theta$, 
of the (standard) Cauchy distribution with density
\[
f(x) = \pi^{-1} (1 + (x-\theta)^2)^{-1}
\]
\begin{enumerate}
%1a)
\item  Using the ``method of scoring'' and the fact that
\[
-\pi^{-1} \int \frac{2x^2-2}{(1+x^2)^3} dx = \frac{1}{2}
\]
find the maximum likelihood estimator for $\theta$ from the sample observations

\vspace{3mm}
\{$-13.71,\,0.40,\,-2.55,\,0.44,\,-0.57,\,-1.05,\,-0.87,\,-1.59,\,-4.11,\,-0.72$\}

\vspace{3mm}

Try to be very explicit about what you do here.

%1b)
\item Interpret the mysterious integral in part (a) in statistical terms.
%1c)
\item Compare the performance of the sample mean and sample median with the 
mle in the Cauchy case.

%1d)
\item Is the mle asymptotically normal in this problem?  
Briefly sketch why or why not.

%1e)
\item Suppose that you wanted to expand the approach in part (a.) so that
instead of a scalar location parameter you had a ``regression effect'', i.e., 
$\theta=x_i'\beta$ for the $i^{\mbox{th}}$ observation.  Describe briefly how you might adapt your 
approach to computing the mle to this case.

\end{enumerate}

%2.
\item Suppose $X_1, \ldots, X_{10}$ constitute a random sample from the density
\[
f(x) = \frac{\theta^3}{2} e^{-\theta x} x^2 \qquad x>0
\]
Knowing only that $\bar X = 1.9, $ find an estimate of $\theta$, and test
the hypothesis that $\theta = 1. $   If possible, try to suggest more than one
test of the hypothesis, compare their conclusions, and discuss their relative
merits from a theoretical standpoint.

%3.
\item In Lecture 18, I tried to connect spline smoothing with model selection 
via the notion of the ``effective dimensionality'' of the fitted function.
In particular, it was suggested that choosing the smoothing parameter 
$\mu=\lambda^{-1} $ by minimizing
\[
\nu(\mu) = n \log (\hat{\sigma} (\mu)) + \frac{1}{2} \log n k(\mu)
\]
where $k(\mu) = Tr (A(\mu))$ might be reasonable.  Suppose several years
from now someone asks you about this.  Try to explain the rationale 
for this expression briefly.  Try to relate it to the application of model
selection in regression.  Suppose someone suggests that instead of using
$k(\mu) = Tr (A(\mu))$ you might want to use $k(u)= n^{-1}Tr (2A(\mu) - 
A^2 (\mu))$, can you suggest a rationale for this?
\end{enumerate}


\vspace{10mm}
\begin{center}
*****
\end{center}

\noindent
University of Illinois \hfill\large{{\bf Official ``Answer'' to Q.3}}\hfill Department of Economics \\
Spring 1997 \hfill \large{{Econ 476 Final Exam 1997}}\hfill Roger Koenker


\vspace{15mm}

This question caused some difficulty in interpretation so I thought it 
might be useful to provide an ``official answer.''  As you will see
there is some rather radical hand waving going on in this answer which 
would be desirable to make more rigorous.

Recall the GCV expression of Craven and Wahba
\[
V(\lambda) = \frac{n^{-1} \parallel (I-A)y\parallel^2}{(n^{-1} Tr (I-A))^2}
\]
The objective is to connect this expression to the proposed SIC expression.
Note first that
\[
\hat \sigma^2 = n^{-1} \parallel (I-A)y\parallel^2
\]
is the usual mle estimate of $\sigma^2$ presuming that $\lambda $ in
$A(\lambda)$ is fixed.  Thus, taking logs we have
\begin{eqnarray*}
\nu(\lambda) &=&  \log V(\lambda) = \log (\hat \sigma^2) - 2 \log (n^{-1} Tr (I-A))\\
&=& \log (\hat \sigma^2 ) -  2 \log (1 - n^{-1} Tr (A))
\end{eqnarray*}
Now using the familiar expansion $\log (1+a) \cong a $ for small $a$ we have,
provided that $Tr(A)$ is small relative to $n$,
\[
\nu(\lambda) \cong \log (\hat\sigma^2) + 2n^{-1} Tr (A)
\]
so minimizing $\nu(\lambda)$ is equivalent to maximizing
\[
\eta (\lambda) = - \frac{n}{2} \log (\hat \sigma^2) - k(\lambda)
\]
where the first term is recognizable as the usual log-likelihood 
evaluated at the mle, for given $\lambda$, and $k(\lambda) = Tr(A)$ is, 
or purports to be, a measure of the dimensionality of the fit.
This is, in effect an interpretation of Craven and Wahba's GCV criterion as 
AIC.  The expression in the exam is simply derived from the above by 
rewriting the likelihood term as $-n\log (\hat\sigma)$ and 
replacing $k(\lambda)$ by $1/2 \log n k(\lambda)$ as in Schwarz.

There is an extensive discussion about alternative definitions of 
$k(\lambda)$ in Hastie and Tibshirani (1990).  One way to motivate the 
suggested alternative to $k(\lambda) = Tr(A(\lambda))$ is to consider 
the usual argument
\begin{eqnarray*}
E\parallel\hat u\parallel^2 &=& E\parallel (I-A)u\parallel^2 \\
&=& \sigma^2 Tr (I-A)^2 \\
&=& \sigma^2 (n-Tr (2A-A^2))
\end{eqnarray*}

So in this approach $Tr (2A-A^2)$ is the dimension of the fitted model.
Of course, if $A$ were really idempotent, then $2A-A^2=A,$ but this 
is not (quite) the case for typical smoothers including the smoothing
spline.


\vspace{10mm} 
\begin{center} 
***** 
\end{center} 

\noindent
University of Illinois \hfill Department of Economics \\
Spring 1995 \hfill \large{{\bf Econ 476 Exam Review}}\hfill Roger Koenker


\vspace{15mm}
\begin{enumerate}
%Q1.
\item Suppose $X_1, \dots , X_{10} $ constitute a random sample from 
the density 
\[
f(x) = \frac{\theta^3}{2} e^{-\theta x} x^2  \qquad \qquad \qquad x>0 .
\]
Knowing only that $\bar{X}=2,$ find an estimate of $\theta $, and test 
the hypothesis that $\theta = 1 .$  Briefly, try to explain any 
special virtues of the procedures you use.

%Q2.
\item In the simple regression through the origin model,
\[
y_i=x_i\beta + u_i
\]
with $x_i$ and $\beta $ scalar.  Explain the significance of the 
condition
\[
\max_i |x_i|^2 / \sum x_i^2 \rightarrow 0
\]
for the asymptotic behavior of $\hat{\beta} = (x'y)/x'x.$

%Q3.
\item Consider the kernel density estimator
\[
\hat{f}_h(x) = n^{-1} \sum K_h (x-X_i) = (nh)^{-1} \sum K_0 \left( 
\frac{x-X_i}{h} \right)
\]
Explain briefly the following argument to compute the asymptotic bias of 
$\hat{f}_n(x)$:
\begin{eqnarray*}
E\hat{f}_h(x) &=& n^{-1}\sum EK_h (x-X_i) \\
 & = & \int K_h (x-u)f(u) du \\
 &=& \int K_0 (s) f(x+sh)ds
\end{eqnarray*}
Thus for $h \rightarrow 0$ we have $E\hat{f}_h(x) \rightarrow f(x) $ and
\begin{eqnarray*}
\mbox{Bias} \: \hat{f}_h(x) &=& \int K_0(s)f(x-sh)ds - f(x) \\
&=& \frac{h^2}{2} \int s^2 K_0(s) ds f '' (x) + o(h^2) 
\end{eqnarray*}
In particular, explain carefully any special assumptions on $K$ or $f$
that are implicitly being used.
%Q4
\item  The celebrated Wald proof of the consistency of the mle is based 
on the following

{\bf Lemma.}  Let $Z_1, ... , Z_n$ be a random sample from the density 
$f(z,\theta_0)$.  Under regularity conditions, on $f$, for any
$\theta \neq \theta_0$, as $n \to \infty$,
\[
(*) \qquad P_{\theta_0} \{ \Pi_{i=1}^n f(Z_i | \theta_0 ) > 
\Pi_{i=1}^n f(Z_i | \theta )\}  \to 1.
\]
The proof may be slightly condensed to the following ``one-liner''
\[
E_{\theta_0} log (f/f_0) < log (E_{\theta_0}  f/f_0 ) = 0.
\]

a.) Explain why $(*)$ implies consistency of the mle.

b.) Briefly describe the regularity conditions.

c.) Explain briefly the role of the regularity conditions.

%\item Explain the following S function.  What does it do and how does it do
%it?
%
%
%{\tt
%lprq\_function(x, y, h, alpha)\{ 
%
%\hspace{10mm} fv <- numeric(length(y)) 
%
%\hspace{10mm} der <- fv
%
%\hspace{10mm} for(i in 1:length(y)) \{ 
%
%\hspace{20mm} z <- x - x[i] 
%
%\hspace{20mm} wx <- dnorm(z/h) 
%
%\hspace {20mm} r <- rq(cbind(1, z) * wx, y * wx, alpha, int = F) 
%
%\hspace{20mm} fv[i] <- r\$coef[1] 
%
%\hspace{20mm} der[i] <- r\$coef[2] 
%
%\hspace{30mm} \} 
%
%\hspace{10mm} res <- y - fv 
%
%\hspace{10mm} list(fv = fv, res = res, der = der) 
%
%\hspace{10mm} \}
%}

%{\tt
%lprq\_function(x, y, h, alpha)\{ \\
%        fv <- numeric(length(y)) \\
%        der <- fv\\
%        for(i in 1:length(y)) \{ \\
%
		%z <- x - x[i] \\
		%wx <- dnorm(z/h) \\
 		%r <- rq(cbind(1, z) * wx, y * wx, alpha, int = F) \\
 		%fv[i] <- r\$coef[1] \\
  		%der[i] <- r\$coef[2]\\ 

%\} \\
%        res <- y - fv \\
%        list(fv = fv, res = res, der = der) \\
%\}
%}
\end{enumerate}



\vspace{10mm} 
\begin{center} 
***** 
\end{center} 

\noindent
University of Illinois \hfill Department of Economics \\
Spring 1996 \hfill \large{{\bf Econ 476 Review}}\hfill Roger Koenker


\vspace{15mm}


\begin{enumerate}
%Q1.
\item Let $X_1, \dots, X_n$ be a random sample from the density 
$f(x, \theta_0)$ and $\hat{\theta}_n$ denote the maximum likelihood estimator of $\theta_o$, 
i.e., $\hat{\theta}_n$ maximizes
\[
\l_n (\theta) = \sum_{i=1}^n \log f(X_i, \theta)
\]
over $\theta \in  \Theta \subseteq \RR^{p}$.  Presuming sufficient regularity 
that $\hat{\theta}_n$ is efficient in the sense that
\[
\sqrt{n}(\hat{\theta}_n-\theta_0) \leadsto \eta (0, I(\theta_0)^{-1})
\]
where $I(\theta_0) $ denotes the Fisher Information of $f$. Sketch the 
argument for fact that 
\[
2(\l_n (\hat{\theta})-\l_n(\theta_0)) \leadsto \chi_{p}^2
\]


%Q2
\item Let $X_1, \dots , X_n$ be a random sample from a df 
$F(x-\theta_0).$  Consider L-estimators of location of the form
\[
\hat{\theta}_0 = \int_0^1 J(u) F_n^{-1}(u)du
\]
where $F_n(\cdot)$ denotes the empirical distribution function constructed 
from the $X_i$'s.  If $F$ has finite Fisher information $I(F)$ and 
a twice continuously differentiable log density, then the optimal $L$-estimator
has weight function of the form
\[
J^*(F(x)) = - \frac{(\log \: f(x))''}{I(F)}
\]
\begin{enumerate}
\item Explain the observation ``note that
\[
\int_{-\infty}^\infty J^*(F(y))dF(y) = \int_0^1 J^*(u)du = 1
\]
and therefore $\hat{\theta}_n $ is location equivariant.'' 

\item The optimality of $\hat{\theta}_n$ may be seen by computing the 
influence function of the general L-estimator as follows:

\begin{description}
\item[\small I.] The $IF$ of the $u^{th}$ sample quantile is 
\[
IF(x, F^{-1}(u), F) = \frac{d}{d\varepsilon} F_\varepsilon^{-1}(u)  = 
\frac{u-\delta_x (F^{-1}(u))}{f(F^{-1}(u))}
\]
which may be shown by differentiating the identity
\[
F_\varepsilon (F_\varepsilon^{-1}(u))=u
\]
where $F_\varepsilon(y) = (1-\varepsilon)F(y) + \varepsilon \delta_x(y)$ to
obtain 
\[
0 = - F(F_\varepsilon^{-1}(u)) + \delta_x (F_\varepsilon^{-1}(y)) + 
f_\varepsilon(F_\varepsilon^{-1}(u))\frac{d}{d\varepsilon}F_\varepsilon^{-1}(u)
\]
and evaluating at $\varepsilon = 0 .$
\item[\small II.] Thus
\begin{eqnarray*}
IF(x, \hat{\theta}_n, F) &=& \int_0^1(J^*(u)(u-\delta_x(F^{-1}(u)))/f(F^{-1}(u))du \\
&=& \int_{-\infty}^\infty J^*(F(y))(F(y)-\delta_x(y))dy\\
&=& \int_{-\infty}^x J^*(F(y))dy - \int_{-\infty}^\infty (1-F(y))J^*(F(y))dy \\
&=& -I(F)^{-1} \int_{-\infty}^x (\log f)'' (y)dy \\
&=& -I(F)^{-1}(\log f)'(x)
\end{eqnarray*}

\item[\small III.] Setting $\psi(x) = - (\log f)'(x) = - f'(x)/f(x)$ we conclude that
$\sqrt{n}(\hat{\theta}_n-\theta_0) \leadsto \eta (0, EIF^2) $ where
\begin{eqnarray*}
EIF^2 &=& \int( \psi^2 (x)/I(F)^2)dF(x) \\
&=& I(F)^{-1}
\end{eqnarray*}
\end{description}
Explain briefly the foregoing result.  Focus on the following aspects
\begin{tabbing}
$\qquad \qquad$ \= (i)(ii)\= How to compute $\theta_n$. \kill
\> (i)\> How to compute $\hat \theta_n$. \\
\> (ii)\> How does $\hat{\theta}_n$ differ from the mle.\\
\> (iii)\> What does the IF tell us about $\hat \theta_n$.
\end{tabbing}
\end{enumerate}


%Q3.
%\item Suppose $X_1, \dots , X_{10} $ constitute a random sample from 
%the density 
%\[
%f(x) = \frac{\theta^3}{2} e^{-\theta x} x^2  \qquad \qquad \qquad x>0 .
%\]
%Knowing only that $\bar{X}=2,$ find an estimate of $\theta $, and test 
%the hypothesis that $\theta = 1 .$  Briefly, try to explain any 
%special virtues of the procedures you use.
%
%Q4
%\item  The celebrated Wald proof of the consistency of the mle is based 
%on the following
%
%{\bf Lemma.}  Let $Z_1, ... , Z_n$ be a random sample from the density 
%$f(z,\theta_0)$.  Under regularity conditions, on $f$, for any
%$\theta \neq \theta_0$, as $n \to \infty$,
%\[
%(*) \qquad P_{\theta_0} \{ \Pi_{i=1}^n f(Z_i | \theta_0 ) > 
%\Pi_{i=1}^n f(Z_i | \theta )\}  \to 1.
%\]
%The proof may be slightly condensed to the following ``one-liner''
%\[
%E_{\theta_0} log (f/f_0) < log (E_{\theta_0}  f/f_0 ) = 0.
%\]
%
%a.) Explain what $(*)$ implies about the consistency of the mle.
%
%b.) Briefly describe the regularity conditions.
%
%c.) Explain the proof in detail.
%
\item Explain the following S function.  What does it do and how does it do
it?


{\tt
lprq\_function(x, y, h, alpha)\{ 

\hspace{10mm} fv <- numeric(length(y)) 

\hspace{10mm} der <- fv

\hspace{10mm} for(i in 1:length(y)) \{ 

\hspace{20mm} z <- x - x[i] 

\hspace{20mm} wx <- dnorm(z/h) 

\hspace {20mm} r <- rq(cbind(1, z) * wx, y * wx, alpha, int = F) 

\hspace{20mm} fv[i] <- r\$coef[1] 

\hspace{20mm} der[i] <- r\$coef[2] 

\hspace{30mm} \} 

\hspace{10mm} res <- y - fv 

\hspace{10mm} list(fv = fv, res = res, der = der) 

\hspace{10mm} \}
}

%{\tt
%lprq\_function(x, y, h, alpha)\{ \\
%        fv <- numeric(length(y)) \\
%        der <- fv\\
%        for(i in 1:length(y)) \{ \\
%
		%z <- x - x[i] \\
		%wx <- dnorm(z/h) \\
 		%r <- rq(cbind(1, z) * wx, y * wx, alpha, int = F) \\
 		%fv[i] <- r\$coef[1] \\
  		%der[i] <- r\$coef[2]\\ 

%\} \\
%        res <- y - fv \\
%        list(fv = fv, res = res, der = der) \\
%\}
%}
\end{enumerate}
\end{document}
